\section{Averaging as a Spectral Approximation}
\label{sec:averaging_spectral}

We now formalize the central claim of this work: under the anisotropic structure
empirically validated in Section~2 and modeled in Section~3, simple averaging
(mean subtraction) provides an effective and conservative approximation to
spectral spike suppression.
In particular, averaging removes a coherent component contained in the spike
subspace while preserving the semantic long tail up to vanishing error.

\subsection{Setup and Notation}

Let $X \in \mathbb{R}^{N\times H}$ denote the matrix of token representations in a
batch, with rows $x_i^\top$.
Recall the decomposition introduced in Section~\ref{sec:model}:
\[
x_i = \ell + s_i,
\qquad
X = \mathbf{1}\ell^\top + S,
\]
where $\ell$ is the shared linguistic component and $S$ collects the semantic
components.

Let $X = U\Sigma V^\top$ be the singular value decomposition of $X$, and let
$P_k = V_{1:k}^\top V_{1:k}$ denote the orthogonal projector onto the top-$k$
right singular vectors (the empirical \emph{spike subspace}).
Define the complementary projector $I-P_k$ as the \emph{tail subspace}.

Let
\[
\mu = \frac{1}{N}\sum_{i=1}^N x_i,
\qquad
X_{\mathrm{avg}} = X - \mathbf{1}\mu^\top
\]
denote the batch mean and the mean-removed matrix.

\subsection{Main Result}

\begin{theorem}[Averaging as conservative spike suppression]
\label{thm:avg_spectral}
Suppose the token representations satisfy Assumptions~1–3 of
Section~\ref{sec:model}, and that the empirical spike subspace $P_k$ contains the
linguistic subspace $\mathcal L$ up to small error (as observed empirically in
Section~2.3).
Then the averaging operator $X_{\mathrm{avg}}$ satisfies the following properties.

\begin{enumerate}

\item \textbf{Mean alignment with the spike subspace.}
The batch mean lies almost entirely in the empirical spike subspace:
\begin{equation}
\label{eq:mean_spike_align}
\frac{\|P_k \mu\|_2^2}{\|\mu\|_2^2} \ge 1 - \varepsilon,
\end{equation}
for modest $k$ and small $\varepsilon$.
Consequently, mean removal suppresses a coherent direction contained in the
spectral spike.
Empirically, this condition holds with $\varepsilon \ll 1$ for $k$ on the order
of the observed spike rank (Section~2.3).

\item \textbf{Spike energy reduction (conservative).}
Averaging removes a nonzero fraction of the spike energy:
\begin{equation}
\label{eq:spike_reduction}
1 - \frac{\|X_{\mathrm{avg}} P_k\|_F^2}{\|X P_k\|_F^2} > 0.
\end{equation}
Because averaging removes a single rank-one component aligned with the batch
mean, the relative spike reduction is largest when the spike subspace is small
and coherent.
As $k$ increases, the same removed direction constitutes a smaller fraction of
the larger spike subspace, and the relative reduction decreases.
This behavior reflects the conservative nature of averaging and is consistent
with the empirical curves in Section~3.

\item \textbf{Tail preservation (additive form).}
Mean removal perturbs the tail component only through the portion of the mean
lying outside the spike subspace:
\begin{equation}
\label{eq:tail_additive_main}
\|X_{\mathrm{avg}}(I-P_k) - X(I-P_k)\|_F
=
\|\mathbf{1}\mu^\top(I-P_k)\|_F
=
\sqrt{N}\,\|(I-P_k)\mu\|_2.
\end{equation}
In particular, when $\mu$ is well captured by the spike subspace (as in
Section~2.3), the additive perturbation to the semantic tail is small.

\smallskip
\noindent\emph{(Relative form under a tail-energy lower bound.)}
If, in addition, the tail energy satisfies $\|X(I-P_k)\|_F^2 \ge cN$ for some
constant $c>0$ (i.e., per-token tail variance is bounded below), then
\begin{equation}
\label{eq:tail_ratio_cor}
\frac{\|X_{\mathrm{avg}}(I-P_k)\|_F^2}{\|X(I-P_k)\|_F^2}
=
1 + O\!\left(\frac{\|(I-P_k)\mu\|_2}{\sqrt{N}}\right),
\end{equation}
so tail energy is preserved up to vanishing relative error as the number of
semantically independent tokens increases.

\end{enumerate}
\end{theorem}

\paragraph{Connection to gradient updates.}
Although the analysis above is stated in terms of token representations, the same
mean–residual separation applies directly to gradient computation in transformer
layers.
In particular, for linear projections, the batch gradient decomposes into a
shared rank-one mean–mean term and a cross-covariance term between error signals
and token-wise variations.
Mean removal eliminates only the shared rank-one component aligned with the
linguistic spike, while preserving the cross-covariance term that carries
token-specific learning signal.
As a result, averaging suppresses harmful coherent structure in gradients without
erasing semantic learning.
A detailed derivation of this gradient decomposition is provided in
Appendix~\ref{app:grad_decomposition}.

\subsection{Interpretation and Comparison with SVD}

Theorem~\ref{thm:avg_spectral} highlights a key distinction between averaging and
SVD-based spike suppression.
Top-$k$ SVD removal eliminates the entire spike subspace and removes increasing
energy as $k$ grows.
In contrast, averaging removes a single coherent direction aligned with the
batch mean.
It therefore provides a conservative approximation: it suppresses the most
harmful shared structure while minimizing collateral removal of semantic
variation.

This distinction explains the empirical curves in Section~3:
averaging removes a substantial fraction of spike energy while preserving almost
all tail energy, whereas SVD-based removal trades off spike suppression against
semantic loss as $k$ increases.

In the next section, we leverage this result to analyze the impact of averaging
on low-bit quantization and adaptive optimization.